% Independently audited by root and four_regular; proof status is recorded in the shared handoff.
% Include after random-hosts-phase6.tex.
\section{Centering the wheel length intervals}
\label{sec:rhost-centered-wheels}

The useful-length construction improves further when each spoke is shortened by approximately one quarter of the rim length.
This centers the long cycle intervals around their required range.
Some outcomes then have invalid spokes; their probability is bounded explicitly below.

\begin{theorem}[An upper coefficient of $199/192$]\label{thm:rhost-centered-wheels}
For the same $p_*$ as in Theorem~\ref{thm:rhost-scale}, and every $p\ge p_*$,
with high probability
\[
 \operatorname{Pex}(G(n,p))
 \le\log_2n+\left(\frac{199}{192}+o(1)\right)\frac{\ln n}{\ln\ln n}.
\]
\end{theorem}

\begin{proof}
Use the distribution from Lemma~\ref{lem:rhost-partial-wheels}, with one change:
put
\[
 B=h-\lfloor R/4\rfloor+s,\qquad w_i=B-X_i.
\]
The total deficit remains $t=h+U$.
Before discarding invalid outcomes, the total edge length is
\[
 4B-t+R=3h+4s-U+R-4\lfloor R/4\rfloor\ge3h+2s.
\]

For a fixed rim arc of length $a$, the corresponding cycle length is uniform on
\[
 [2B-t+a,2B+a]\cap\mathbb Z.
\]
The probability that this length lies outside $[h+1,2h]$ is at most
\[
 \frac{|a-R/2|+O(s)}h.
\]
The constant in this estimate is absolute, uniformly over all choices.
For each of the two unit-rim pairs, the two arcs are $1,R-1$.
Their two absolute deviations sum to $R-2$.
For the other four pairs, both arcs differ from $R/2$ by $O(s)$.
Thus, if $k$ counts useful non-rim labels,
\[
 \mathbb E(12-k)\le\frac{2\mathbb E R+O(s)}h=\frac1{12}+o(1),
 \qquad \mathbb E(k-3)\ge\frac{107}{12}-o(1).
\]

We bound the probability $p_{\rm bad}$ that a spoke is nonpositive.
The marginal mass formula is
\[
 \Pr(X_i=t-r)=\frac{c_{t-r}(2r+1)c_r}{t+1}.
\]
Stirling's formula gives $c_j\sim(\pi j)^{-1/2}$.
Riemann sums, together with $c_j\le C/\sqrt{j+1}$ at the endpoints, show that $X_i/h$
converges in distribution to a variable with density
\[
 f(y)=\frac2\pi y^{-1/2}(1-y)^{1/2},\qquad 0<y<1.
\]
Also, $R/h$ converges to an independent uniform variable on $[0,1/12]$.
The independence follows because $R$ is independent of $U$ and the deficit distribution.
The limiting event $X_i/h\ge1-R/(4h)$ has a boundary of probability zero.
A union bound therefore gives
\[
 \limsup p_{\rm bad}
 \le 12\int_0^{1/12}\frac8\pi
       \int_0^{x/4}\frac{\sqrt u}{\sqrt{1-u}}\,du\,dx.
\]
For $0\le x\le1/12$,
\[
 \frac8\pi\int_0^{x/4}\frac{\sqrt u}{\sqrt{1-u}}\,du
 \le\frac{2x^{3/2}}{3\pi\sqrt{1-x/4}}
 \le\frac9{40}x^{3/2}.
\]
For the last inequality, use $\pi\ge3$ and $\sqrt{47/48}>80/81$.
Consequently,
\[
 p_{\rm bad}\le\frac9{100}\left(\frac1{12}\right)^{3/2}+o(1)
             =\frac3{800\sqrt3}+o(1).
\]

Discard outcomes with a nonpositive spoke or repeated non-rim cycle lengths.
The latter event still has probability $o(1)$.
Discarding a bad-spoke outcome loses at most nine in value.
Discard outcomes with $k\le3$, whose values are nonpositive.
The retained expected value is at least
\[
 \frac{107}{12}-\frac{27}{800\sqrt3}-o(1).
\]
The maximum weighted degree remains $(12+o(1))/h$.
Indeed, the pair-deficit interval still has $h+U+1$ equally likely positions, and the rim distribution is unchanged.
The weighted codegree remains $O(h^{-3/2})$.
The shared-index split proof is unchanged.
For complementary pairs, condition on $R,L$ and use the independent variable $U$.
For two arcs of one pair, use the same $R,L$ smoothing as before.
For a rim and a long length, fixing $R$ and then a pair sum costs $O(h^{-2})$.
Thus Lemma~\ref{lem:rhost-weighted-matching} supplies disjoint marked sets of total value at least
\[
 \left(\frac{107}{144}-\frac9{3200\sqrt3}-o(1)\right)h.
\]

Every retained metric has the required local density.
Its total length is at least $3h$ and its largest arm is at most $h+s$.
A three-spoke rank-three core therefore has length at least $2h$.
A four-spoke rank-three core omits at most one rim arm, of length at most $R/2+O(s)$, and also exceeds $2h$.
For a two-spoke rank-two core, include the full rim and use
\[
 2B-t+R\ge h+R/2.
\]
A rank-two core with at least three spokes has length at least $2h-R+s>h$.
Every non-rim cycle has length at least $h-R/2\gg\sqrt h$, while the rim has length at least $6s$.
The edge count is $O(h)$ and the same four-edge deletion gives a spanning path.
The wheel embedding estimate therefore applies.

Replace baseline single cycles by these marked wheel sets, as in Theorem~\ref{thm:rhost-one-fifty-one}.
The short-block rank is at most
\[
 \left(\frac{149}{144}+\frac9{3200\sqrt3}+o(1)\right)h.
\]
Since $\sqrt3>5/3$ and $27/16000<1/576$, the constant is less than $199/192$.
The previous exposure scheme and excess calculation finish the proof.
\end{proof}

The wheel checker also verifies 3000 centered metrics against all 42 kernel core types.
These finite checks support the metric formulas; the distributional limit and random matching arguments are proved above.
